%%%%%%%%%%  scan idx 183  |  folio 170  %%%%%%%%%%
$F\,|\,S_{f_{ij}}$ est
engendré par un nombre fini de sections $h_{ijk}$ au-dessus de $S_{f_{ij}}$. Il existe (EGA I 9.3.1 et IV 1.7.5) un entier $q_i\ge0$ tel que $h_{ijk}\otimes f_{ij}^{\otimes q_i}$ se prolonge en une section de $F\otimes L_i^{\otimes q_im_i}$ au-dessus de $S$. Autrement dit, $F\otimes L_i^{\otimes q_im_i}$ est engendré au-dessus de $S_i$ par un nombre fini de ses sections au-dessus de $S$. Procédant comme dans la démonstration de (EGA II 4.5.5), on en déduit qu'il existe un entier $n_i$, qu'on peut supposer $>0$, tel que, pour tout $n\ge n_i$, $F\otimes L_i^{\otimes n}$ soit engendré au-dessus de $S_i$ par un nombre fini de ses sections au-dessus de $S$. En particulier, il existe un entier $r_i\ge0$ et un morphisme
\[
\mathcal O_S^{r_i}\otimes L_i^{\otimes-n_i}\longrightarrow F
\]
qui induit un épimorphisme au-dessus de $S_i$. L'implication (i$'$)$\Rightarrow$(ii) en résulte aussitôt.

(ii)$\Rightarrow$(iii) : Comme tout faisceau quasi-cohérent sur $S$ est limite inductive de ses sous-faisceaux quasi-cohérents de type fini (EGA I 9.4.9 et IV 1.7.7), on peut se borner à prouver (iii) pour $F$ de type fini. Il existe alors, pour $i\in I$ des morphismes
\[
u_i:\mathcal O_S^{r_i}\otimes L_i^{\otimes-n_i}\longrightarrow F
\]
tels que $\sum u_i$ soit un épimorphisme. Or on a
\[
\begin{aligned}
\operatorname{Im}(u_i)&=\operatorname{Im}\bigl(u_i\otimes\operatorname{Id}(L_i^{\otimes n_i})\otimes\operatorname{Id}(L_i^{\otimes-n_i})\bigr)\\
&=\operatorname{Im}\bigl(u_i\otimes\operatorname{Id}(L_i^{\otimes n_i})\bigr)\otimes L_i^{\otimes-n_i}
\end{aligned}
\]
Mais $\operatorname{Im}(u_i\otimes\operatorname{Id}(L_i^{\otimes n_i}))$ est un sous-faisceau de $F_{(i,n_i)}$, d'où l'assertion.

(iii)$\Rightarrow$(iii$'$) : C'est trivial.

(iii$'$)$\Rightarrow$(i) : Soient $x\in S$, $U$ un ouvert contenant $x$, $F$ un Idéal 
quasi-%
